% Original pages: 89-96
% Figures: none
% Uncertainty log: none

\section{Noetherian Rings}

We recall that a ring $A$ is said to be \emph{Noetherian} if it satisfies the following three
equivalent conditions:

1) Every non-empty set of ideals in $A$ has a maximal element.

2) Every ascending chain of ideals in $A$ is stationary.

3) Every ideal in $A$ is finitely generated.

(The equivalence of these conditions was proved in \amref{6.1} and \amref{6.2}.)

Noetherian rings are by far the most important class of rings in commutative
algebra: we have seen some examples already in Chapter 6.  In this chapter we
shall first show that Noetherian rings reproduce themselves under various
familiar operations---in particular we prove the famous basis theorem of Hilbert.
We then proceed to make a number of important deductions from the
Noetherian condition, including the existence of primary decompositions.

\medskip
\noindent\textit{\textbf{Proposition 7.1.\amtarget{7.1}} If $A$ is Noetherian and $\phi$ is a homomorphism of $A$ onto a
ring $B$, then $B$ is Noetherian.}

\noindent\emph{Proof.}  This follows from \amref{6.6}, since $B\cong A/\mathfrak a$, where $\mathfrak a=\operatorname{Ker}(\phi)$.\proofbox

\medskip
\noindent\textit{\textbf{Proposition 7.2.\amtarget{7.2}} Let $A$ be a subring of $B$; suppose that $A$ is Noetherian and
that $B$ is finitely generated as an $A$-module.  Then $B$ is Noetherian (as a ring).}

\noindent\emph{Proof.}  By \amref{6.5} $B$ is Noetherian as an $A$-module, hence also as a $B$-module.\proofbox

\noindent\textbf{Example.}  $B=\mathbb Z[\ii]$, the ring of Gaussian integers.  By \amref{7.2} $B$ is Noetherian.
More generally, the ring of integers in any algebraic number field is Noetherian.

\medskip
\noindent\textit{\textbf{Proposition 7.3.\amtarget{7.3}} If $A$ is Noetherian and $S$ is any multiplicatively closed
subset of $A$, then $S^{-1}A$ is Noetherian.}

\noindent\emph{Proof.}  By (3.11---i) and (1.17---iii) the ideals of $S^{-1}A$ are in one-to-one order-preserving correspondence with the contracted ideals of $A$, hence satisfy the maximal condition.  (Alternative proof: if $\mathfrak a$ is any ideal of $A$, then $\mathfrak a$ has a finite set
of generators, say $x_1,\ldots,x_n$, and it is clear that $S^{-1}\mathfrak a$ is generated by $x_1/1,\ldots,
x_n/1$.)\proofbox

\medskip
\noindent\textit{\textbf{Corollary 7.4.\amtarget{7.4}} If $A$ is Noetherian and $\mathfrak p$ is a prime ideal of $A$, then $A_\mathfrak p$ is
Noetherian.}\proofbox

\medskip
\noindent\textit{\textbf{Theorem 7.5.\amtarget{7.5}} (Hilbert's Basis Theorem). If $A$ is Noetherian, then the
polynomial ring $A[x]$ is Noetherian.}

\noindent\emph{Proof.}  Let $\mathfrak a$ be an ideal in $A[x]$.  The leading coefficients of the polynomials in
$\mathfrak a$ form an ideal $I$ in $A$.  Since $A$ is Noetherian, $I$ is finitely generated, say by
$a_1,\ldots,a_n$.  For each $i=1,\ldots,n$ there is a polynomial $f_i\in A[x]$ of the form
$f_i=a_ix^{r_i}+$ (lower terms).  Let $r=\max_{i=1}^n r_i$.  The $f_i$ generate an ideal
$\mathfrak a'\subseteq\mathfrak a$ in $A[x]$.

Let $f=ax^m+$ (lower terms) be any element of $\mathfrak a$; we have $a\in I$.  If $m\geq r$,
write $a=\sum_{i=1}^n u_i a_i$, where $u_i\in A$; then $f-\sum u_if_ix^{m-r_i}$ is in $\mathfrak a$ and has degree
$<m$.  Proceeding in this way, we can go on subtracting elements of $\mathfrak a'$ from $f$
until we get a polynomial $g$, say, of degree $<r$; that is, we have $f=g+h$,
where $h\in\mathfrak a'$.

Let $M$ be the $A$-module generated by $1,x,\ldots,x^{r-1}$; then what we have
proved is that $\mathfrak a=(\mathfrak a\cap M)+\mathfrak a'$.  Now $M$ is a finitely generated $A$-module,
hence is Noetherian by \amref{6.5}, hence $\mathfrak a\cap M$ is finitely generated (as an $A$-module)
by \amref{6.2}.  If $g_1,\ldots,g_m$ generate $\mathfrak a\cap M$ it is clear that the $f_i$ and the $g_j$ generate $\mathfrak a$.
Hence $\mathfrak a$ is finitely generated and so $A[x]$ is Noetherian.\proofbox

\emph{Remark.}  It is also true that $A$ Noetherian $\Rightarrow A[[x]]$ Noetherian ($A[[x]]$ being
the ring of formal power series in $x$ with coefficients in $A$).  The proof runs
almost parallel to that of \amref{7.5} except that one starts with the terms of lowest
degree in the power series belonging to $\mathfrak a$.  See also \amref{10.27}.

\medskip
\noindent\textit{\textbf{Corollary 7.6.\amtarget{7.6}} If $A$ is Noetherian so is $A[x_1,\ldots,x_n]$.}

\noindent\emph{Proof.}  By induction on $n$ from \amref{7.5}.\proofbox

\medskip
\noindent\textit{\textbf{Corollary 7.7.\amtarget{7.7}} Let $B$ be a finitely-generated $A$-algebra.  If $A$ is Noetherian,
then so is $B$.}

\noindent\textit{
In particular, every finitely-generated ring, and every finitely generated
algebra over a field, is Noetherian.}

\noindent\emph{Proof.}  $B$ is a homomorphic image of a polynomial ring $A[x_1,\ldots,x_n]$, which is
Noetherian by \amref{7.6}.\proofbox

\medskip
\noindent\textit{\textbf{Proposition 7.8.\amtarget{7.8}} Let $A\subseteq B\subseteq C$ be rings.  Suppose that $A$ is Noetherian,
that $C$ is finitely generated as an $A$-algebra and that $C$ is either (i) finitely
generated as a $B$-module or (ii) integral over $B$.  Then $B$ is finitely generated
as an $A$-algebra.}

\noindent\emph{Proof.}  It follows from \amref{5.1} and \amref{5.2} that the conditions (i) and (ii) are equivalent in this situation.  So we may concentrate on (i).

Let $x_1,\ldots,x_m$ generate $C$ as an $A$-algebra, and let $y_1,\ldots,y_n$ generate $C$ as a
$B$-module.  Then there exist expressions of the form
\begin{equation}
x_i=\sum_j b_{ij}y_j\qquad(b_{ij}\in B)
\tag{1}
\end{equation}
\begin{equation}
y_iy_j=\sum_k b_{ijk}y_k\qquad(b_{ijk}\in B).
\tag{2}
\end{equation}
Let $B_0$ be the algebra generated over $A$ by the $b_{ij}$ and the $b_{ijk}$.  Since $A$ is
Noetherian, so is $B_0$ by \amref{7.7}, and $A\subseteq B_0\subseteq B$.

Any element of $C$ is a polynomial in the $x_i$ with coefficients in $A$.  Substituting (1) and making repeated use of (2) shows that each element of $C$ is a
linear combination of the $y_j$ with coefficients in $B_0$, and hence $C$ is finitely
generated as a $B_0$-module.  Since $B_0$ is Noetherian, and $B$ is a submodule of $C$,
it follows (by \amref{6.5} and \amref{6.2}) that $B$ is finitely generated as a $B_0$-module.  Since
$B_0$ is finitely generated as an $A$-algebra, it follows that $B$ is finitely-generated as
an $A$-algebra.\proofbox

\medskip
\noindent\textit{\textbf{Proposition 7.9.\amtarget{7.9}} Let $k$ be a field, $E$ a finitely generated $k$-algebra.  If $E$ is a
field then it is a finite algebraic extension of $k$.}

\noindent\emph{Proof.}  Let $E=k[x_1,\ldots,x_n]$.  If $E$ is not algebraic over $k$ then we can renumber the $x_i$ so that $x_1,\ldots,x_r$ are algebraically independent over $k$, where
$r\geq1$, and each of $x_{r+1},\ldots,x_n$ is algebraic over the field $F=k(x_1,\ldots,x_r)$.
Hence $E$ is a finite algebraic extension of $F$ and therefore finitely generated as an
$F$-module.  Applying \amref{7.8} to $k\subseteq F\subseteq E$, it follows that $F$ is a finitely generated
$k$-algebra, say $F=k[y_1,\ldots,y_s]$.  Each $y_j$ is of the form $f_j/g_j$, where $f_j$ and $g_j$
are polynomials in $x_1,\ldots,x_r$.

Now there are infinitely many irreducible polynomials in the ring
$k[x_1,\ldots,x_r]$ (adapt Euclid's proof of the existence of infinitely many prime
numbers).  Hence there is an irreducible polynomial $h$ which is prime to each
of the $g_j$ (for example, $h=g_1g_2\cdots g_s+1$ would do) and the element $h^{-1}$ of $F$
could not be a polynomial in the $y_j$.  This is a contradiction.  Hence $E$ is algebraic over $k$, and therefore finite algebraic.\proofbox

\medskip
\noindent\textit{\textbf{Corollary 7.10.\amtarget{7.10}} Let $k$ be a field, $A$ a finitely generated $k$-algebra.  Let $\mathfrak m$ be a
maximal ideal of $A$.  Then the field $A/\mathfrak m$ is a finite algebraic extension of $k$.
In particular, if $k$ is algebraically closed then $A/\mathfrak m\cong k$.}

\noindent\emph{Proof.}  Take $E=A/\mathfrak m$ in \amref{7.9}.\proofbox

\amref{7.10} is the so-called ``weak'' version of Hilbert's Nullstellensatz
($={}$ theorem of the zeros).  The proof given here is due to Artin and Tate.  For
its geometrical meaning, and the ``strong'' form of the theorem, see the Exercises
at the end of this chapter.

\subsection{Primary decomposition in Noetherian rings}

The next two lemmas show that every ideal $\ne(1)$ in a Noetherian ring has a
primary decomposition.

An ideal $\mathfrak a$ is said to be \emph{irreducible} if
\[
\mathfrak a=\mathfrak b\cap\mathfrak c\Rightarrow(\mathfrak a=\mathfrak b\text{ or }\mathfrak a=\mathfrak c).
\]

\medskip
\noindent\textit{\textbf{Lemma 7.11.\amtarget{7.11}} In a Noetherian ring $A$ every ideal is a finite intersection of
irreducible ideals.}

\noindent\emph{Proof.}  Suppose not; then the set of ideals in $A$ for which the lemma is false is
not empty, hence has a maximal element $\mathfrak a$.  Since $\mathfrak a$ is reducible, we have
$\mathfrak a=\mathfrak b\cap\mathfrak c$ where $\mathfrak b\supset\mathfrak a$ and $\mathfrak c\supset\mathfrak a$.  Hence each of $\mathfrak b,\mathfrak c$ is a finite intersection of
irreducible ideals and therefore so is $\mathfrak a$: contradiction.\proofbox

\medskip
\noindent\textit{\textbf{Lemma 7.12.\amtarget{7.12}} In a Noetherian ring every irreducible ideal is primary.}

\noindent\emph{Proof.}  By passing to the quotient ring, it is enough to show that if the zero ideal
is irreducible then it is primary.  Let $xy=0$ with $y\ne0$, and consider the chain
of ideals $\operatorname{Ann}(x)\subseteq\operatorname{Ann}(x^2)\subseteq\cdots$.  By the a.c.c., this chain is stationary, i.e., we
have $\operatorname{Ann}(x^n)=\operatorname{Ann}(x^{n+1})=\cdots$ for some $n$.  It follows that $(x^n)\cap(y)=0$;
for if $a\in(y)$ then $ax=0$, and if $a\in(x^n)$ then $a=bx^n$, hence $bx^{n+1}=0$, hence
$b\in\operatorname{Ann}(x^{n+1})=\operatorname{Ann}(x^n)$, hence $bx^n=0$; that is, $a=0$.  Since $(0)$ is irreducible
and $(y)\ne0$ we must therefore have $x^n=0$, and this shows that $(0)$ is primary.\proofbox

From these two lemmas we have at once

\medskip
\noindent\textit{\textbf{Theorem 7.13.\amtarget{7.13}} In a Noetherian ring $A$ every ideal has a primary decomposition.}\proofbox

Hence all the results of Chapter 4 apply to Noetherian rings.

\medskip
\noindent\textit{\textbf{Proposition 7.14.\amtarget{7.14}} In a Noetherian ring $A$, every ideal $\mathfrak a$ contains a power of
its radical.}

\noindent\emph{Proof.}  Let $x_1,\ldots,x_k$ generate $r(\mathfrak a)$: say $x_i^{n_i}\in\mathfrak a$ $(1\leq i\leq k)$.  Let $m=
\sum_{i=1}^k(n_i-1)+1$.  Then $r(\mathfrak a)^m$ is generated by the products $x_1^{r_1}\cdots x_k^{r_k}$ with
$\sum r_i=m$; from the definition of $m$ we must have $r_i\geq n_i$ for at least one index $i$,
hence each such monomial lies in $\mathfrak a$, and therefore $r(\mathfrak a)^m\subseteq\mathfrak a$.\proofbox

\medskip
\noindent\textit{\textbf{Corollary 7.15.\amtarget{7.15}} In a Noetherian ring the nilradical is nilpotent.}

\noindent\emph{Proof.}  Take $\mathfrak a=(0)$ in \amref{7.14}.\proofbox

\medskip
\noindent\textit{\textbf{Corollary 7.16.\amtarget{7.16}} Let $A$ be a Noetherian ring, $\mathfrak m$ a maximal ideal of $A$, $\mathfrak q$ any
ideal of $A$.  Then the following are equivalent:}

\noindent\textit{i) $\mathfrak q$ is $\mathfrak m$-primary;}

\noindent\textit{ii) $r(\mathfrak q)=\mathfrak m$;}

\noindent\textit{iii) $\mathfrak m^n\subseteq\mathfrak q\subseteq\mathfrak m$ for some $n>0$.}

\noindent\emph{Proof.}  i) $\Rightarrow$ ii) is clear; ii) $\Rightarrow$ i) from \amref{4.2}; ii) $\Rightarrow$ iii) from \amref{7.14}; iii) $\Rightarrow$ ii) by
taking radicals: $\mathfrak m=r(\mathfrak m^n)\subseteq r(\mathfrak q)\subseteq r(\mathfrak m)=\mathfrak m$.\proofbox

\medskip
\noindent\textit{\textbf{Proposition 7.17.\amtarget{7.17}} Let $\mathfrak a\ne(1)$ be an ideal in a Noetherian ring.  Then the
prime ideals which belong to $\mathfrak a$ are precisely the prime ideals which occur in the
set of ideals $(\mathfrak a:x)$ $(x\in A)$.}

\noindent\emph{Proof.}  By passing to $A/\mathfrak a$ we may assume that $\mathfrak a=0$.  Let $\bigcap_{i=1}^n\mathfrak q_i=0$ be a
minimal primary decomposition of the zero ideal, and let $\mathfrak p_i$ be the radical of $\mathfrak q_i$.

Let $\mathfrak a_i=\bigcap_{j\ne i}\mathfrak q_j\ne0$.  Then from the proof of \amref{4.5} we have $r(\operatorname{Ann}(x))=\mathfrak p_i$
for any $x\ne0$ in $\mathfrak a_i$, so that $\operatorname{Ann}(x)\subseteq\mathfrak p_i$.

Since $\mathfrak q_i$ is $\mathfrak p_i$-primary, by \amref{7.14} there exists an integer $m$ such that $\mathfrak p_i^m\subseteq\mathfrak q_i$,
and therefore $\mathfrak a_i\mathfrak p_i^m\subseteq\mathfrak a_i\cap\mathfrak p_i^m\subseteq\mathfrak a_i\cap\mathfrak q_i=0$.  Let $m\geq1$ be the smallest
integer such that $\mathfrak a_i\mathfrak p_i^m=0$, and let $x$ be a non-zero element in $\mathfrak a_i\mathfrak p_i^{m-1}$.  Then
$\mathfrak p_ix=0$, therefore for such an $x$ we have $\operatorname{Ann}(x)\supseteq\mathfrak p_i$, and hence $\operatorname{Ann}(x)=\mathfrak p_i$.

Conversely, if $\operatorname{Ann}(x)$ is a prime ideal $\mathfrak p$, then $r(\operatorname{Ann}(x))=\mathfrak p$ and so by \amref{4.5}
$\mathfrak p$ is a prime ideal belonging to 0.\proofbox

\subsection{Exercises}

\noindent\textbf{1.} Let $A$ be a non-Noetherian ring and let $\Sigma$ be the set of ideals in $A$ which are not
finitely generated.  Show that $\Sigma$ has maximal elements and that the maximal
elements of $\Sigma$ are prime ideals.

[Let $\mathfrak a$ be a maximal element of $\Sigma$, and suppose that there exist $x,y\in A$ such that
$x\notin\mathfrak a$ and $y\notin\mathfrak a$ and $xy\in\mathfrak a$.  Show that there exists a finitely generated ideal
$\mathfrak a_0\subseteq\mathfrak a$ such that $\mathfrak a_0+(x)=\mathfrak a+(x)$, and that $\mathfrak a=\mathfrak a_0+x\cdot(\mathfrak a:x)$.  Since
$(\mathfrak a:x)$ strictly contains $\mathfrak a$, it is finitely generated and therefore so is $\mathfrak a$.]

Hence a ring in which every prime ideal is finitely generated is Noetherian
(I. S. Cohen).

\medskip
\noindent\textbf{2.} Let $A$ be a Noetherian ring and let $f=\sum_{n=0}^\infty a_nx^n\in A[[x]]$.  Prove that $f$ is
nilpotent if and only if each $a_n$ is nilpotent.

\medskip
\noindent\textbf{3.} Let $\mathfrak a$ be an irreducible ideal in a ring $A$.  Then the following are equivalent:

i) $\mathfrak a$ is primary;

ii) for every multiplicatively closed subset $S$ of $A$ we have $(S^{-1}\mathfrak a)^c=(\mathfrak a:x)$ for
some $x\in S$;

iii) the sequence $(\mathfrak a:x^n)$ is stationary, for every $x\in A$.

\medskip
\noindent\textbf{4.} Which of the following rings are Noetherian?

i) The ring of rational functions of $z$ having no pole on the circle $|z|=1$.

ii) The ring of power series in $z$ with a positive radius of convergence.

iii) The ring of power series in $z$ with an infinite radius of convergence.

iv) The ring of polynomials in $z$ whose first $k$ derivatives vanish at the origin
($k$ being a fixed integer).

v) The ring of polynomials in $z,w$ all of whose partial derivatives with respect
to $w$ vanish for $z=0$.

In all cases the coefficients are complex numbers.

\medskip
\noindent\amstar\textbf{5.} Let $A$ be a Noetherian ring, $B$ a finitely generated $A$-algebra, $G$ a finite group of
$A$-automorphisms of $B$, and $B^G$ the set of all elements of $B$ which are left fixed
by every element of $G$.  Show that $B^G$ is a finitely generated $A$-algebra.

\medskip
\noindent\textbf{6.} If a finitely generated ring $K$ is a field, it is a finite field.

[If $K$ has characteristic 0, we have $\mathbb Z\subset\mathbb Q\subseteq K$.  Since $K$ is finitely generated
over $\mathbb Z$ it is finitely generated over $\mathbb Q$, hence by \amref{7.9} is a finitely generated $\mathbb Q$-module.  Now apply \amref{7.8} to obtain a contradiction.  Hence $K$ is of characteristic
$p>0$, hence is finitely generated as a $\mathbb Z/(p)$-algebra.  Use \amref{7.9} to complete the
proof.]

\medskip
\noindent\textbf{7.} Let $X$ be an affine algebraic variety given by a family of equations $f_\alpha(t_1,\ldots,t_n)
=0$ ($\alpha\in I$) (Chapter 1, Exercise 27).  Show that there exists a finite subset $I_0$ of $I$
such that $X$ is given by the equations $f_\alpha(t_1,\ldots,t_n)=0$ for $\alpha\in I_0$.

\medskip
\noindent\textbf{8.} If $A[x]$ is Noetherian, is $A$ necessarily Noetherian?

\medskip
\noindent\textbf{9.} Let $A$ be a ring such that

(1) for each maximal ideal $\mathfrak m$ of $A$, the local ring $A_\mathfrak m$ is Noetherian;

(2) for each $x\ne0$ in $A$, the set of maximal ideals of $A$ which contain $x$ is
finite.

Show that $A$ is Noetherian.

[Let $\mathfrak a\ne0$ be an ideal in $A$.  Let $\mathfrak m_1,\ldots,\mathfrak m_r$ be the maximal ideals which
contain $\mathfrak a$.  Choose $x_0\ne0$ in $\mathfrak a$ and let $\mathfrak m_1,\ldots,\mathfrak m_{r+s}$ be the maximal ideals
which contain $x_0$.  Since $\mathfrak m_{r+1},\ldots,\mathfrak m_{r+s}$ do not contain $\mathfrak a$ there exist $x_j\in\mathfrak a$ such
that $x_j\notin\mathfrak m_{r+j}$ $(1\leq j\leq s)$.  Since each $A_{\mathfrak m_i}$ $(1\leq i\leq r)$ is Noetherian, the extension of $\mathfrak a$ in $A_{\mathfrak m_i}$ is finitely generated.  Hence there exist $x_{s+1},\ldots,x_t$ in $\mathfrak a$
whose images in $A_{\mathfrak m_i}$ generate $A_{\mathfrak m_i}\mathfrak a$ for $i=1,\ldots,r$.  Let $\mathfrak a_0=(x_0,\ldots,x_t)$.
Show that $\mathfrak a_0$ and $\mathfrak a$ have the same extension in $A_\mathfrak m$ for every maximal ideal $\mathfrak m$, and
deduce by \amref{3.9} that $\mathfrak a_0=\mathfrak a$.]

\medskip
\noindent\textbf{10.} Let $M$ be a Noetherian $A$-module.  Show that $M[x]$ (Chapter 2, Exercise 6) is a
Noetherian $A[x]$-module.

\medskip
\noindent\textbf{11.} Let $A$ be a ring such that each local ring $A_\mathfrak p$ is Noetherian.  Is $A$ necessarily
Noetherian?

\medskip
\noindent\textbf{12.} Let $A$ be a ring and $B$ a faithfully flat $A$-algebra (Chapter 3, Exercise 16).  If $B$
is Noetherian, show that $A$ is Noetherian.  [Use the ascending chain condition.]

\medskip
\noindent\textbf{13.} Let $f:A\longrightarrow B$ be a ring homomorphism of finite type and let $f^*:\operatorname{Spec}(B)\longrightarrow
\operatorname{Spec}(A)$ be the mapping associated with $f$.  Show that the fibers of $f^*$ are
Noetherian subspaces of $B$.

\medskip
\noindent\emph{Nullstellensatz, strong form}

\medskip
\noindent\textbf{14.} Let $k$ be an algebraically closed field, let $A$ denote the polynomial ring
$k[t_1,\ldots,t_n]$ and let $\mathfrak a$ be an ideal in $A$.  Let $V$ be the variety in $k^n$ defined by the
ideal $\mathfrak a$, so that $V$ is the set of all $x=(x_1,\ldots,x_n)\in k^n$ such that $f(x)=0$ for all
$f\in\mathfrak a$.  Let $I(V)$ be the ideal of $V$, i.e. the ideal of all polynomials $g\in A$ such that
$g(x)=0$ for all $x\in V$.  Then $I(V)=r(\mathfrak a)$.

[It is clear that $r(\mathfrak a)\subseteq I(V)$.  Conversely, let $f\notin r(\mathfrak a)$, then there is a prime ideal $\mathfrak p$
containing $\mathfrak a$ such that $f\notin\mathfrak p$.  Let $\bar f$ be the image of $f$ in $B=A/\mathfrak p$, let $C=B_{\bar f}=
B[1/\bar f]$, and let $\mathfrak m$ be a maximal ideal of $C$.  Since $C$ is a finitely generated $k$-algebra we have $C/\mathfrak m\cong k$, by \amref{7.9}.  The images $x_i$ in $C/\mathfrak m$ of the generators $t_i$
of $A$ thus define a point $x=(x_1,\ldots,x_n)\in k^n$, and the construction shows that
$x\in V$ and $f(x)\ne0$.]

\medskip
\noindent\amstar\textbf{15.} Let $A$ be a Noetherian local ring, $\mathfrak m$ its maximal ideal and $k$ its residue field, and
let $M$ be a finitely generated $A$-module.  Then the following are equivalent:

i) $M$ is free;

ii) $M$ is flat;

iii) the mapping of $\mathfrak m\otimes M$ into $A\otimes M$ is injective;

iv) $\operatorname{Tor}_1^A(k,M)=0$.

[To show that iv) $\Rightarrow$ i), let $x_1,\ldots,x_n$ be elements of $M$ whose images in $M/\mathfrak mM$
form a $k$-basis of this vector space.  By \amref{2.8}, the $x_i$ generate $M$.  Let $F$ be a free
$A$-module with basis $e_1,\ldots,e_n$ and define $\phi:F\longrightarrow M$ by $\phi(e_i)=x_i$.  Let $E=\operatorname{Ker}
(\phi)$.  Then the exact sequence $0\longrightarrow E\longrightarrow F\longrightarrow M\longrightarrow0$ gives us an exact sequence
\[
0\longrightarrow k\otimes_A E\longrightarrow k\otimes_A F\xrightarrow{1\otimes\phi}k\otimes_A M\longrightarrow0.
\]
Since $k\otimes F$ and $k\otimes M$ are vector spaces of the same dimension over $k$, it
follows that $1\otimes\phi$ is an isomorphism, hence $k\otimes E=0$, hence $E=0$ by
Nakayama's Lemma ($E$ is finitely generated because it is a submodule of $F$, and $A$
is Noetherian).]

\medskip
\noindent\amstar\textbf{16.} Let $A$ be a Noetherian ring, $M$ a finitely generated $A$-module.  Then the following
are equivalent:

i) $M$ is a flat $A$-module;

ii) $M_\mathfrak p$ is a free $A_\mathfrak p$-module, for all prime ideals $\mathfrak p$;

iii) $M_\mathfrak m$ is a free $A_\mathfrak m$-module, for all maximal ideals $\mathfrak m$.

In other words, flat $={}$ locally free.  [Use Exercise 15.]

\medskip
\noindent\textbf{17.} Let $A$ be a ring and $M$ a Noetherian $A$-module.  Show (by imitating the proofs
of \amref{7.11} and \amref{7.12}) that every submodule $N$ of $M$ has a primary decomposition
(Chapter 4, Exercises 20--23).

\medskip
\noindent\textbf{18.} Let $A$ be a Noetherian ring, $\mathfrak p$ a prime ideal of $A$, and $M$ a finitely generated
$A$-module.  Show that the following are equivalent:

i) $\mathfrak p$ belongs to 0 in $M$;

ii) there exists $x\in M$ such that $\operatorname{Ann}(x)=\mathfrak p$;

iii) there exists a submodule of $M$ isomorphic to $A/\mathfrak p$.

Deduce that there exists a chain of submodules
\[
0=M_0\subset M_1\subset\cdots\subset M_r=M
\]
such that each quotient $M_i/M_{i-1}$ is of the form $A/\mathfrak p_i$, where $\mathfrak p_i$ is a prime ideal
of $A$.

\medskip
\noindent\textbf{19.} Let $\mathfrak a$ be an ideal in a Noetherian ring $A$.  Let
\[
\mathfrak a=\bigcap_{i=1}^r\mathfrak b_i=\bigcap_{j=1}^s\mathfrak c_j
\]
be two minimal decompositions of $\mathfrak a$ as intersections of \emph{irreducible} ideals.  Prove
that $r=s$ and that (possible after re-indexing the $\mathfrak c_j$) $r(\mathfrak b_i)=r(\mathfrak c_i)$ for all $i$.

[Show that for each $i=1,\ldots,r$ there exists $j$ such that
\[
\mathfrak a=\mathfrak b_1\cap\cdots\cap\mathfrak b_{i-1}\cap\mathfrak c_j\cap\mathfrak b_{i+1}\cap\cdots\cap\mathfrak b_r.
\]
]

State and prove an analogous result for modules.

\medskip
\noindent\textbf{20.} Let $X$ be a topological space and let $\mathcal F$ be the smallest collection of subsets of $X$
which contains all open subsets of $X$ and is closed with respect to the formation
of finite intersections and complements.

i) Show that a subset $E$ of $X$ belongs to $\mathcal F$ if and only if $E$ is a finite union of sets
of the form $U\cap C$, where $U$ is open and $C$ is closed.

ii) Suppose that $X$ is irreducible and let $E\in\mathcal F$.  Show that $E$ is dense in $X$
(i.e., that $\bar E=X$) if and only if $E$ contains a non-empty open set in $X$.

\medskip
\noindent\textbf{21.} Let $X$ be a Noetherian topological space (Chapter 6, Exercise 5) and let $E\subseteq X$.
Show that $E\in\mathcal F$ if and only if, for each irreducible closed set $X_0\subseteq X$, either
$\overline{E\cap X_0}\ne X_0$ or else $E\cap X_0$ contains a non-empty open subset of $X_0$.  [Suppose
$E\notin\mathcal F$.  Then the collection of closed sets $X'\subseteq X$ such that $E\cap X'\notin\mathcal F$ is not
empty and therefore has a minimal element $X_0$.  Show that $X_0$ is irreducible and
then that each of the alternatives above leads to the conclusion that $E\cap X_0\in\mathcal F$.]
The sets belonging to $\mathcal F$ are called the \emph{constructible subsets} of $X$.

\medskip
\noindent\textbf{22.} Let $X$ be a Noetherian topological space and let $E$ be a subset of $X$.  Show that $E$
is open in $X$ if and only if, for each irreducible closed subset $X_0$ in $X$, either
$E\cap X_0=\varnothing$ or else $E\cap X_0$ contains a non-empty open subset of $X_0$.  [The
proof is similar to that of Exercise 21.]

\medskip
\noindent\textbf{23.} Let $A$ be a Noetherian ring, $f:A\longrightarrow B$ a ring homomorphism of finite type (so
that $B$ is Noetherian).  Let $X=\operatorname{Spec}(A)$, $Y=\operatorname{Spec}(B)$ and let $f^*:Y\longrightarrow X$ be
the mapping associated with $f$.  Then the image under $f^*$ of a constructible
subset $E$ of $Y$ is a constructible subset of $X$.

[By Exercise 20 it is enough to take $E=U\cap C$ where $U$ is open and $C$ is closed
in $Y$; then, replacing $B$ by a homomorphic image, we reduce to the case where $E$
is open in $Y$.  Since $Y$ is Noetherian, $E$ is quasi-compact and therefore a finite
union of open sets of the form $\operatorname{Spec}(B_g)$.  Hence reduce to the case $E=Y$.  To
show that $f^*(Y)$ is constructible, use the criterion of Exercise 21.  Let $X_0$ be an
irreducible closed subset of $X$ such that $f^*(Y)\cap X_0$ is dense in $X_0$.  We have
$f^*(Y)\cap X_0=f^*(f^{*-1}(X_0))$, and $f^{*-1}(X_0)=\operatorname{Spec}((A/\mathfrak p)\otimes_A B)$, where $X_0=
\operatorname{Spec}(A/\mathfrak p)$.  Hence reduce to the case where $A$ is an integral domain and $f$ is injective.  If $Y_1,\ldots,Y_n$ are the irreducible components of $Y$, it is enough to show that
some $f^*(Y_i)$ contains a non-empty open set in $X$.  So finally we are brought down
to the situation in which $A,B$ are integral domains and $f$ is injective (and still
of finite type); now use Chapter 5, Exercise 21 to complete the proof.]

\medskip
\noindent\textbf{24.} With the notation and hypotheses of Exercise 23, $f^*$ is an open mapping $\Longleftrightarrow$
$f$ has the going-down property (Chapter 5, Exercise 10).  [Suppose $f$ has the
going-down property.  As in Exercise 23, reduce to proving that $E=f^*(Y)$ is
open in $X$.  The going-down property asserts that if $\mathfrak p\in E$ and $\mathfrak p'\subseteq\mathfrak p$, then
$\mathfrak p'\in E$: in other words, that if $X_0$ is an irreducible closed subset of $X$ and $X_0$
meets $E$, then $E\cap X_0$ is dense in $X_0$.  By Exercises 20 and 22, $E$ is open in $X$.]

\medskip
\noindent\textbf{25.} Let $A$ be Noetherian, $f:A\longrightarrow B$ of finite type and \emph{flat} (i.e., $B$ is flat as an $A$-module).  Then $f^*:\operatorname{Spec}(B)\longrightarrow\operatorname{Spec}(A)$ is an open mapping.  [Exercise 24 and
Chapter 5, Exercise 11.]
